\chapter{Being grammatical isn't always enough}

\section{ideas}
\begin{enumerate}
    \item Haspelmath's (1999) notion of `extravagance' as one of the keys to language change, we account for the systematic patterning of deliberate linguistic subversion by appealing to tension between the need to stand out and the need to remain intelligible.
    \item Zwicky \& Pullum (1987), ``Plain Morphology and Expressive Morphology,''
\end{enumerate}

\textit{I have 30 years.} I don't put an asterisk there because it's fully grammatical. If somebody asked you how long you have until you can retire, you might, at a certain age, say, \textit{I have 30 years}. But if somebody asked you for your age, \textit{I have 30 years} is no good as a response. It's just not how we say it in English. Instead, we say \textit{I'm 30 years old}. Clearly, we \textbf{could} say it that way. That's how they do it in French (\textit{J'ai 30 ans.}) and many other languages. And, if, in response to a question about age, a French speaker told you in English \textit{I have 30 years}, you'd almost certainly understand. But for some reason, that formulation is just weird in English.

Conversely, if you want to say, ``press 2 on your phone'' in French, you would be mistaken to say, \textit{presser 2}. Instead, you would use \textit{faites le 2}, literally `do the 2'. Again, I don't use an asterisk because \textit{presser 2} is perfectly grammatical French in other contexts; it's just not what the French say for dialing a phone, and they would not accept it as correct.

\bigskip

Yesterday my wife texted me \textit{you still at the gym?} to which I replied jokingly \textit{I still at the gym}. Any humour in my reply relies on a superficial similarity to her question paired with a shared understanding that \textit{I still at the gym} is ungrammatical. But it seems to me that it goes beyond that. It also relies subtly on the understanding that the \textit{you still at the gym?} construction is only legitimate in a casual conversation.

\ea
    \ea[]{\textit{Is the device dependable?}}
    \ex[?]{\textit{the device dependable?}}
    \z
\z
\ea
    \ea[]{\textit{Is it possible to consider this issue globally?}}
    \ex[*]{\textit{It possible to consider this issue globally?}}
    \z
\z
\ea
    \ea[]{\textit{Is this cause for complaint?}}
    \ex[?]{\textit{This cause for complaint?}}
    \z
\z
\ea
    \ea[]{\textit{Is she any the poorer for never having the chance to live amid those extirpated species}}
    \ex[?]{\textit{She any the poorer for never having the chance to live amid those extirpated species}}
    \z
\z
\ea
    \ea[]{\textit{Are the pages of a book a place?}}
    \ex[?]{\textit{The pages of a book a place?}}
    \z
\z
\ea
    \ea[]{\textit{Are the authors saying that we are free to read any way we wish?}}
    \ex[?]{\textit{The authors saying that we are free to read any way we wish?}}
    \z
\z
\ea
    \ea[]{\textit{Are the FEASP-strategies used in daily instruction?}}
    \ex[*]{\textit{The FEASP-strategies used in daily instruction?}}
    \z
\z

\section{Russian dolls} \label{sec:dolls}

A \textsc{relative clause} is a clause that typically functions as a modifier in a noun phrase. For example, in \textit{I found the key that I had dropped}, there is a noun phrase, \textit{the key that I had dropped}, and \textit{that I had dropped} is a relative clause functioning as a modifier in that noun phrase. In English, a relative clause is what linguists call an \textsc{unbounded dependency construction}. First, I'll explain the \textsc{dependency} part, and then I'll tell you why it's (probably) unbounded.

Instead of writing \textit{I found the key that I had dropped} as one sentence, we could split it up into two: \textit{I found the key} and \textit{I had dropped it}. By doing so, the dependency relationship becomes a little easier to see: the meaning of \textit{it} in the second clause depends on \textit{key} in the first. If the first clause were \textit{I found the money}, then \textit{it} would have to mean `money' and not `key'. When we write the original sentence, with the relative clause, the \textit{it} goes away. But there's still a kind of dependency relationship because \textit{I had dropped} means something on its own, but what it means is `I, myself, had fallen'. To get the right interpretation, that I had dropped the key and not my body, we need a dependency relationship between \textit{dropped} and \textit{key}. That's the dependency part.

When linguists say this is \textit{unbounded}, what they mean is that, in principle at least, there's no limit to how far away \textit{key} and \textit{dropped} are. There's no boundary limit on how many clauses deep the dependency can go.
\begin{itemize}
    \item \textit{I found the key that I dropped (it).}
    \item \textit{I found the key that I had dropped (it).}
    \item \textit{I found the key that I thought I had dropped (it).}
    \item \textit{I found the key that I thought she said I had dropped (it).}
    \item \textit{I found the key that I thought she said you heard I had dropped (it).}
    \item Etc.
\end{itemize}

David Gil, the linguist who studies Riau Indonesian, claims that, in that language, relative clauses are grammatical. But he also claims that speakers of Riau Indonesian reject constructions in which the dependency relationship crosses a clause boundary. In other words, a translation of \textit{I found the key I dropped} with a relative clause is possible, but \textit{I found the key that I thought I had dropped} is not. 

As with \textit{I have 30 years} or \textit{do the 2}, you would most likely be understood, but you would not be speaking Riau Indonesian in a way that is accepted by the Riau Indonesians. This time, though, it's not just a particular idiom for expressing one's age or instructing someone on how to navigate a phone menu; it's a whole class of sentences that are rejected. You can't say, \textit{this is the rat that at the cheese that sat in the house that Jack built}. You can't say, \textit{I know that you know that it's wrong}. And you can't say, \textit{It would be good if you just tried the one I gave you}. 

It's not that these things can't be expressed -- though there \textit{are} propositions that are simply not entertained in some languages. It's just that a competent speaker of Riau Indonesian would just put the clauses one after the other instead of one inside the other. English speakers do this too. There's the example about the key from above, but there's also the somewhat informal and rather threatening \textit{touch this and you die}.

Dan Everett, a linguist who spent many years in the Amazon learning and describing a language called Pirahã, would regularly run into situations where he was told, ``we don't talk about that.'' This was not a case of ``we don't say it that way,'' but literally, that's not a topic that we discuss. Such topics are not taboo, like sex or money might be, they're simply things that don't interest them, such as the number of items in a group, or history or~\dots

\section{Conventions}

It's perfectly true that some ways of expressing things are conventional, and some are not. But sometimes convention breaking leads to turns of phrase that that are creative, poetic, artistic, and liable to delight an audience. Other times, convention breaking language is just a little odd. What's the difference.

One of the most well-known of artful convention breakers was William Shakespeare, but whether you like Shakespeare or not, you probably don't want to hear, yet again, about his brilliance with formulating new figures of speech, such as \textit{star-crossed lovers}, \textit{wild-goose chase}, and \textit{break the ice}. So, I won't bore you with cases like \textit{green-eyed monster}, \textit{in a pickle}, and \textit{bide one's time}.

And you may not be a fan of hip-hop, so I won't bring up the linguistic creativity that has given rise to expressions such as \textit{bling bling} (B.G.), \textit{cray} (Jay-Z), and \textit{fomo} (Missy Elliott).

And I certainly won't mention all the new idioms such as \textit{embiggen}, \textit{meh}, and \textit{worst day ever} that have been birthed by the writers of The Simpsons.

Instead, I'll introduce you to Anne Carson, a contemporary Canadian poet, essayist, and classicist. While not as widely recognized as some literary giants, her blending of prose and poetry in works like \textit{Nox} showcases a unique linguistic creativity that challenges traditional literary boundaries.

\subsection{Syntactic Satiation}


\section{Information structure isn't syntax}

Back in the 1970s, when bell-bottoms were groovy and disco was king, linguists had some funny ideas about grammar. They didn't know much about how context affects language use -- what we now call discourse conditions. So when they noticed that some sentences sounded weird, they chalked it up to hard-and-fast syntactic rules.

Take this pair of sentences:

\begin{quote}
\textit{Jack says that he'll take out the recycling.}\\
\textsuperscript{?}\textit{Jack says that the recycling will be taken out by him.}
\end{quote}

The second one sounds off, doesn't it? But why? 

In those days, linguists would slap an asterisk on that second sentence faster than you could say ``Stayin' Alive.'' They thought it was a syntactic no-no. Their rule went something like this: ``the noun phrase in a passive by-phrase can't refer to the subject of the main clause.''

But here's where it gets interesting. My friend and co-author Geoff Pullum recently shared with me an example that turns this old idea on its head. It's from Richard Osman's novel \textit{We Solve Murders} (which, by the way, is a cracking good read):

\begin{quote}
\textit{Henk shoos Trouble away, which only makes Trouble more determined to be stroked by him.}
\end{quote}

Now, if you're scratching your head wondering who Henk and Trouble are, don't worry. All you need to know is that Henk is a cat-hating Dutchman, and Trouble is a very persistent feline.

According to those 1970s rules, this sentence should be a grammatical disaster. We've got a passive by-phrase (\textit{by him}) referring back to the subject (Henk). It's breaking all the rules! And yet... it works. It sounds perfectly fine. What gives?

The answer lies in something called \textsc{information packaging}. It's not about syntax at all, but about how we present information in a conversation. The general rule is that the noun phrase in a passive by-phrase should refer to something at least as new in the discourse as the subject.

That's why ``\textit{Darling, I love you}'' sounds great, but ``\textsuperscript{?}\textit{Darling, you are loved by me}'' sounds like you're auditioning for a b-grade romance novel. The ``\textit{by me}'' part is old news -- we already know the sentence is about the speaker.

So why does Osman's sentence work? Because the alternatives are even worse:

\begin{quote}
\textsuperscript{??}\textit{Henk shoos Trouble away, which only makes Trouble more determined to have him stroke him.}
\end{quote}

Or how about:

\begin{quote}
\textsuperscript{?}\textit{Henk shoos Trouble away, which only makes Trouble more determined that Henk should stroke him.}
\end{quote}

Better, but still clunky. Osman's version, breaking the ``rule'', actually communicates more clearly than either of these alternatives.

\section{Reasons}
It may not be the case that all of these share any one common element apart from that they've caught on. In this section, I'll suggest some elements that I think are likely to increase or decrease our chances of accepting a particular locution.

    \subsection{Group membership}
        \subsubsection*{Intransitive verbs} \label{sec:intrans}
            
            There are certain verbs in English that we think of as transitive. In other words, we expect those verbs to have a relationship to a noun phrase object. An example is the \textit{dropped} sentence we considered in Section \ref{sec:dolls}. We can say \textit{I dropped the keys} or \textit{the keys that I dropped}, or even \textit{the keys were dropped}, and in all three cases, \textit{dropped} has a relationship with \textit{the keys}. Verbs with such relationships are called \textsc{transitive verbs}. If, instead, I just said, \textit{I dropped}, meaning `I fell to the floor', without one of these relationships between \textit{dropped} and a noun phrase, then we call it an \textsc{intransitive verb}.
            
            Some verbs, like \textit{drop} are usually transitive, and some like \textit{exist} never are. Many verbs are happy to be in a transitive relationship, but they're also perfectly content by themselves. Verbs like this include \textit{eat}, \textit{play}, \textit{move}, \textit{start}, \textit{turn}, etc.
            
            But within certain communities, verbs that we expect to be transitive may become intransitive. For instance, many English speakers would find it odd to be asked, \textit{what do you think about lifting?} The response might be \textit{lifting what?} because the person is expecting \textit{lift} to be a transitive verb. If you're a weightlifter, though, \textit{lift} is intransitive, and the question is about lifting weights. Other verbs like this are \textit{ride} (a bike), \textit{drink} (alcohol), \textit{use} (drugs), [additional examples needed].
            
            \textit{Brew} is widely associated with brewing beer (though we also brew coffee, tea, [additional examples needed]). Nevertheless, if a co-worker told you \textit{I've stopped brewing}, as you waited for them to finish with the photocopier, you might find the construction a bit odd, and you'd probably be inclined to ask what it was they had been brewing, at least if you wanted the conversation to develop. Beer might not even occur to you. In contrast, if the same co-worker told you \textit{I've stopped drinking}, it would be quite clear that it's some kind of alcohol. And this would be true even if you had no previous knowledge about your co-worker's level of alcohol consumption. In the community of beer-making hobbyists, though, \textit{I've stopped brewing} would probably be just as clear as \textit{I've stopped drinking}.
    
            Many verbs can, in principle, be \textit{intransitive}. Which is to say that there's nothing in the grammar that prohibits it, as far as we can tell. But membership in 

        \subsubsection*{Passives and dangling modifiers}
    
            The converse of this is when certain groups prohibit or at least frown on constructions that are really quite common. Among people who consider themselves writers, it's quite common to find those who object strongly to sentences like \textit{It was produced in Canada}.
            
            A clause like this is said to be a \textsc{passive clause}. Passive clauses take whatever it is that typically follows the verb -- \textit{lift weights}, for example, \textit{brew beer}, or \textit{produce it} -- and makes it the subject: \textit{weights are lifted}, \textit{beer is brewed}, and \textit{it is produced in Canada}. 
            
            Jeff Brown, a friend of mine, was doing a PhD in philosophy under a passive hating professor (who went on to be his thesis supervisor). He realized this when he received the following feedback on a piece of writing: ``Avoid passive voice -- and don't write short sentences. Look at Nietzsche.''

            Setting aside that Jeff was reading Nietzsche in translation, we checked his copy of the book they were studying. There was a passive in the second paragraph.

            \begin{center}
                -- --
            \end{center}
            
            Read virtually any great novel, and you will come across \textit{passive clauses}. If you record people speaking with friends, at weddings, in formal debates, while watching the game, or really in any possible setting at any level of formality, people use passives, and nobody blinks an eye. There are passives in Shakespeare and passives in \textit{The Simpsons} and passives in \textit{The New York Times}. You'll find passives in advertising, passives in pop songs, and passives in children's books. They are ubiquitous. In normal English, passive constructions appear frequently, though the exact frequency would depend on the context and the corpus analyzed. But writers, at least when they're writing or teaching or editing, seem to believe that 
            \textst{passives should be avoided} (ahem) one should avoid passives. The claim is not that they are ungrammatical, not at all, but rather that they are insidious, mendacious, corrupting, and just downright bad, bad, bad, bad, bad.

            There's no evidence of anyone having had any problem with the passive until George Orwell wrote ``Politics and the English language.'' [is this true]? [etc]

            \begin{center}
                -- --
            \end{center}
            
            It often takes training to identify \textsc{dangling modifiers}, while most grammatical mistakes are obvious to any competent speaker of the language. Second, even those who know how to spot danglers don't have the intuitive feeling we typically experience when the grammar is violated.
            
            For this reason, among others [is that really why?], Jim Donaldson, at the University of Edinburgh, has argued that danglers are fully grammatical. Why, then, can they be so bad? Donaldson argues that some danglers, those he calls ``howlers'', are bad because they're unclear. This is the same kind of badness we experience when a novel has a character dropping a gun in one sentence and drawing it in the next. The reader is brought up short, and will often go back and re-read. Writing about drawing a gun that isn't there clearly isn't grammatical, but it is certainly confusing. Howlers, it appears, are like that.



\section{Pragmatics} 
\begin{quote}
    [Pragmatics] also provides a necessary foundation for explaining such mysteries as how it is that words can change their meanings over time (Traugott \& Dasher 2002) and how it is that children are able to learn a language in the first place (Tomasello 2003; Goldberg 2006).\cite[17]{Israel2011}
\end{quote}